\normaltrue \difficilefalse \tdifficilefalse
\correctionfalse

%\UPSTIidClasse{11} % 11 sup, 12 spé
%\newcommand{\UPSTIidClasse}{11}

\exer{Fourchette}
\label{G2:01:2003}
\setcounter{question}{0}\marginnote{\xpComp{PPM}{02}}%\UPSTIcompetence[2]{G2-01}
\index{Compétence G2-01}\index{Compétence PPM-02}

\ifcorrection
\else
\marginnote{\textbf{Pas de corrigé pour cet exercice.}}
\fi


\ifprof 
\else
Soit la pièce suivante.
\begin{center}
\includegraphics[width=.45\linewidth]{fig_01}
\includegraphics[width=.45\linewidth]{fig_02}
\end{center}
 \fi
 
\question{Proposer une gamme de fabrication, depuis la préparation du brut jusqu'à la finition.}
\ifprof

\else 
\fi

Avant usinage de la pièce, une simulation a été réalisée afin d'observer les contraintes et les déplacements engendrés dans la pièce.

\begin{marginfigure}
\includegraphics[width=\linewidth]{fig_02}
\end{marginfigure}

\question{La pièce est réalisée en acier. Définir un acier, puis définir une fonte.}

\question{La limite élastique vaut 800 MPa. Expliquer la signification de cette valeur et détailler le protocole expérimental permettant de la déterminer.}


\ifprof
\else

\marginnote{Corrigé voir \ref{G2:01:2003}.}

\fi

% BEGIN AUTO CORRIGE INCLUDE
\InclureCorrige{corrige.tex}
% END AUTO CORRIGE INCLUDE
